\documentclass[11pt]{article}
\usepackage[a4paper,margin=27mm]{geometry}
\usepackage{amsmath,amssymb,amsthm,mathtools,booktabs}
\usepackage{microtype}
\usepackage[hidelinks]{hyperref}
\emergencystretch=2em
\newtheorem{theorem}{Theorem}[section]
\newtheorem{lemma}[theorem]{Lemma}
\newtheorem{proposition}[theorem]{Proposition}
\newtheorem{corollary}[theorem]{Corollary}
\theoremstyle{definition}\newtheorem{definition}[theorem]{Definition}
\theoremstyle{remark}\newtheorem{remark}[theorem]{Remark}
\newcommand{\T}{\mathbb T}
\title{A sharp universal profile for multivariate\\ characteristic-function factorization}
\author{Alexandr Martinevski\\
\small Independent Researcher, Lithuania\\
\small \href{mailto:alex.martinevski@gmail.com}{\texttt{alex.martinevski@gmail.com}}\\
\small ORCID: \href{https://orcid.org/0009-0000-4230-4414}{0009-0000-4230-4414}}
\date{September 2026}
\begin{document}
\maketitle

\begin{abstract}
For complex random variables $Z_1,\ldots,Z_d$ with $|Z_j|\le1$, prescribe the moduli
$r_j=|\mathbb EZ_j|$, put $\theta_j=\arccos r_j$, $R=\prod_jr_j$ and $s=\sum_j\theta_j$. We prove the exact profile
\[
\sup\left\{\left|\mathbb E\prod_jZ_j-\prod_j\mathbb EZ_j\right|:|Z_j|\le1,\ |\mathbb EZ_j|=r_j\right\}
=R-\cos(\min\{s,\pi\}).
\]
The sharp universal characteristic-function factorization constant is therefore the one-dimensional corollary
\[
C_d=\max_{0\le x\le\pi}\{\cos^d(x/d)-\cos x\}
=\max_{0\le\theta\le\pi/d}\{\cos^d\theta-\cos(d\theta)\}
=\|T_d-x^d\|_{\infty,[-1,1]}.
\]
It is attained by a common fair sign and is simultaneously sharp for mutual-independence, sub-independence and
one-variable characteristic-function defects. The profile yields a short proof based only on Cauchy--Schwarz and
log-concavity of cosine. We also record selected consequences of the real-parameter interpolation
$C(p)=\max_{0\le\theta\le\pi/p}\{\cos^p\theta-\cos(p\theta)\}$: $C(p)$ is strictly increasing for $p>2$, every maximizer lies in
$(\pi/(p+1),\pi/p)$, and the classical point $\pi/(p+1)$ is asymptotically sharp. In particular,
\[
C_d-\left\{\cos\frac\pi{d+1}+\cos^d\frac\pi{d+1}\right\}
=\frac{\pi^6}{8d^4}+O(d^{-5}),
\]
which gives a four-term expansion of $C_d$. We retain exact real-valued and Gaussian comparison results and an
algebraic description of the critical point.
\end{abstract}

\section{Introduction}
Mutual independence of random blocks $X_j\in\mathbb R^{p_j}$ is equivalent to the factorization
$\varphi_X(t_1,\dots,t_d)=\prod_j\varphi_{X_j}(t_j)$. Empirical-characteristic-function tests of total independence go
back at least to \cite{Csorgo1985,Kankainen1995}; related $L^2$ constructions include distance multivariance and kernel
criteria \cite{BKRS2019,Pfister2018}. For two blocks, Daniu\v{s}is et al.\ introduced the uniform Fourier dependence
measure (UFDM), the $L^\infty$ norm of the joint-minus-product characteristic function, and proved the universal bound
one \cite{Daniusis2025}. The weaker diagonal factorization property is sub-independence
\cite{EHSV2010,Hamedani2013,Schennach2019}.

The present line of work has an earlier origin. In an unpublished student research project of the author (1996),
the uniform characteristic-function factorization defect was introduced as a dependence quantity, with particular
motivation from cryptography; the limiting value $2$ and Gaussian calculations were already part of that project.
The exact finite-$d$ profile, the sharp constants and all proofs in the present paper are new developments from 2026.

The first question here is not estimation but the population extremal problem: how large can the uniform
factorization defect be for $d$ blocks? A previous formulation of this paper answered this by a direct angular proof.
The stronger result below gives the entire sharp profile after the marginal mean moduli are fixed. The universal
constant $C_d$ then follows in a few lines. This reorganization has two advantages: the proof is shorter, and the
geometry of the extremizer is visible before the final one-dimensional maximization.

Our second aim is to record only those consequences of the continuous interpolation
\[
C(p)=\max_{0\le\theta\le\pi/p}\{\cos^p\theta-\cos(p\theta)\},\qquad p>2,
\]
that sharpen the original sequence $C_d=C(d)$. The broader two-parameter cosine/Chebyshev surface and its scaling
limits are left for separate work. In particular, we use the interpolation to obtain continuous monotonicity and a
sharp asymptotic correction to the elementary point $\pi/(d+1)$.

The fixed-marginal problem is treated separately in \cite{MartinevskiMarginals}; in particular, that companion work gives exact results for several symmetric unimodal marginals, a general large-$d$ ``price of marginals'', and a non-symmetric extension. The local arithmetic, irreducibility and Galois theory of the critical polynomial introduced below are developed in \cite{MartinevskiCritical}.

The identity $C_d=\|T_d-x^d\|_\infty$ should not be confused with the classical Chebyshev minimax problem for monic
polynomials, where the optimally scaled polynomial is $2^{1-d}T_d$ \cite{Rivlin1990}. Our searches did not locate the
profile theorem below, the sharp $d$-linear complex unit-ball constant, or the sequence $C_d$ for $d\ge3$; this is a
statement about the searches performed, not a claim of priority. Abstract Gr\"uss and multiplicativity-defect
inequalities form relevant background, but they do not by themselves give the constant below.

\section{Exact fixed-mean profile and the universal constant}
For $r=(r_1,\ldots,r_d)\in[0,1]^d$ define
\[
\mathcal C_d(r_1,\ldots,r_d)=\sup\left\{\left|\mathbb E\prod_{j=1}^dZ_j-\prod_{j=1}^d\mathbb EZ_j\right|:
|Z_j|\le1,\ |\mathbb EZ_j|=r_j\right\}.
\]
Put
\[
\theta_j=\arccos r_j,\qquad R=\prod_{j=1}^dr_j,\qquad s=\sum_{j=1}^d\theta_j.
\]

\begin{theorem}[Exact profile]\label{thm:profile}
For every $d\ge1$ and $r\in[0,1]^d$,
\[
\boxed{\mathcal C_d(r_1,\ldots,r_d)=R-\cos(\min\{s,\pi\}).}
\]
If $s\le\pi$, equality is attained by $Z_j=e^{i\varepsilon\theta_j}$ with a common fair sign
$\mathbb P(\varepsilon=\pm1)=1/2$. If $s\ge\pi$, equality is attained by a coupling with
$\prod_jZ_j\equiv-1$ and the prescribed means.
\end{theorem}

For random blocks $X=(X_1,\ldots,X_d)$ define
\[
\Delta_X(t)=\varphi_X(t_1,\ldots,t_d)-\prod_{j=1}^d\varphi_{X_j}(t_j),\qquad
D_d(X)=\sup_t|\Delta_X(t)|,
\]
and let $D_d^\star$ be the supremum over all joint laws and all block dimensions.

\begin{corollary}[Sharp universal constant]\label{thm:main}
For every $d\ge2$,
\[
\boxed{
D_d^\star=C_d:=\max_{0\le x\le\pi}\{\cos^d(x/d)-\cos x\}
=\max_{0\le\theta\le\pi/d}\{\cos^d\theta-\cos(d\theta)\}
=\|T_d-x^d\|_{L^\infty[-1,1]}.
}
\]
Equivalently,
\[
\left|\mathbb E\prod_jZ_j-\prod_j\mathbb EZ_j\right|\le C_d\qquad(|Z_j|\le1),
\]
and the bound is sharp. It is attained by $Z_j=e^{i\varepsilon\theta_d^*}$, or by
$X_1=\cdots=X_d=\varepsilon$ at the common frequency $\theta_d^*$.
\end{corollary}

\begin{corollary}\label{cor:forms}
The same constant is sharp for the sub-independence defect
\[
D_d^{\rm sub}(X)=\sup_t\left|\varphi_{X_1+\cdots+X_d}(t)-\prod_j\varphi_{X_j}(t)\right|,
\]
for the one-variable defect $S_d(X)=\sup_t|\varphi_X(dt)-\varphi_X(t)^d|$, and after restricting $D_d$ to exchangeable
laws or to laws with identical marginals.
\end{corollary}

\begin{theorem}\label{thm:unique}
For $d\ge3$, $g_d(\theta)=\cos^d\theta-\cos(d\theta)$ has a unique maximizer $\theta_d^*\in(0,\pi/d)$. With
$c_d^*=\cos\theta_d^*$,
\[
U_{d-1}(c_d^*)=(c_d^*)^{d-1}.
\]
The maximizers of $|T_d(x)-x^d|$ on $[-1,1]$ are exactly $\pm c_d^*$. Moreover $C_2<C_3<C_4<\cdots$ and $C_d\uparrow2$.
\end{theorem}

\section{Proof of the profile theorem and Corollary~\ref{thm:main}}
\begin{proof}[Proof of Theorem~\ref{thm:profile}]
Multiplying coordinates by deterministic unimodular constants, assume $\mathbb EZ_j=r_j\ge0$. We argue by induction on
$d$; the case $d=1$ is trivial. Put
\[
W=\prod_{j<d}Z_j,\quad u=\mathbb EW,\quad R'=\prod_{j<d}r_j,\quad
s'=\sum_{j<d}\theta_j,\quad r=r_d=\cos\theta,\quad \theta=\theta_d.
\]
Then
\[
\mathbb E(WZ_d)-rR'=\mathbb E[(W-u)(Z_d-r)]+r(u-R'),
\]
and Cauchy--Schwarz gives
\begin{equation}\label{eq:csprofile}
|\Delta|\le \sin\theta\sqrt{1-|u|^2}+r\,t,\qquad t:=|u-R'|,
\end{equation}
because $\mathbb E|Z_d-r|^2\le1-r^2$ and $\mathbb E|W-u|^2\le1-|u|^2$.

By induction,
\[
0\le t\le R'-\cos(\min\{s',\pi\}).
\]
The reverse triangle inequality gives $|u|\ge|R'-t|$. Write
$R'-t=\cos\phi$ with
\[
\phi\in[\arccos R',\min\{s',\pi\}].
\]
Then $\sqrt{1-|u|^2}\le\sin\phi$, and \eqref{eq:csprofile} yields
\[
|\Delta|\le \sin\theta\sin\phi+\cos\theta(R'-\cos\phi)
=R-\cos(\theta+\phi).
\]
If $s\le\pi$, then $\theta+\phi\le s$ and the right-hand side is at most $R-\cos s$. If $s\ge\pi$, the admissible
interval contains $\phi=\pi-\theta$ (since $\arccos R'\le\pi/2\le\pi-\theta$ and $s'\ge\pi-\theta$), so the upper bound is
$R+1$. This proves the required upper bound.

If $s\le\pi$, take a common fair sign $\varepsilon$ and $Z_j=e^{i\varepsilon\theta_j}$. Then
$\mathbb EZ_j=r_j$ and $\mathbb E\prod_jZ_j=\cos s$, giving $R-\cos s$.

Suppose $s\ge\pi$ and set
\[
Q=\left\{r\in[0,1]^d:\sum_j\arccos r_j\ge\pi\right\}.
\]
The set $Q$ is compact and convex because $\arccos$ is concave. Every extreme point is either on the boundary
$\sum_j\arccos r_j=\pi$ or is a cube vertex: if the sum is strictly larger than $\pi$ and some coordinate lies in
$(0,1)$, a sufficiently small perturbation of that coordinate in both directions stays in $Q$. A cube vertex belongs to
$Q$ exactly when it has at least two zero coordinates. Boundary points are realized by the common-sign construction
with $\sum_j\theta_j=\pi$, for which $\prod_jZ_j=-1$. A vertex with $k\ge2$ zero coordinates is realized by taking the
remaining coordinates equal to $1$, choosing $U_1,\ldots,U_{k-1}$ independently uniform on the unit circle, and setting
$U_k=-\overline{U_1\cdots U_{k-1}}$; then every zero-coordinate mean is $0$ and the product is identically $-1$.
By the finite-dimensional Minkowski theorem, every $r\in Q$ is a convex combination of extreme points; mixing the
corresponding couplings preserves the means and keeps the product identically $-1$. Hence the value $1+R$ is attained.
\end{proof}

\begin{proof}[Proof of Corollary~\ref{thm:main}]
Let $\theta_j=\arccos r_j$. If $s=\sum_j\theta_j\le\pi$, strict concavity of $\log\cos$ gives
\[
R=\prod_j\cos\theta_j\le\cos^d(s/d),
\]
with equality when the angles are equal. Thus Theorem~\ref{thm:profile} gives
\[
\mathcal C_d(r)\le\cos^d(s/d)-\cos s.
\]
If $s\ge\pi$, the same concavity and monotonicity give
$R\le\cos^d(s/d)\le\cos^d(\pi/d)$, so
$\mathcal C_d(r)=1+R\le1+\cos^d(\pi/d)$, the endpoint value at $s=\pi$. Therefore
\[
\sup_r\mathcal C_d(r)=\max_{0\le s\le\pi}\{\cos^d(s/d)-\cos s\}=C_d.
\]
Equal angles and the common fair sign attain every value in this one-dimensional family, hence in particular its
maximum. Taking $Z_j=e^{i\langle t_j,X_j\rangle}$ proves the characteristic-function bound, while scalar angle-valued
variables realize the extremizer.

Finally, for $x\in[-1,1]$, parity reduces $|T_d(x)-x^d|$ to $x\ge0$; writing $x=\cos\theta$ gives the defect of the
common-sign construction, so it is at most $C_d$. Equality occurs at $x=\cos\theta_d^*$.
\end{proof}

\begin{proof}[Proof of Corollary~\ref{cor:forms}]
Each restricted problem is bounded by the full one. For a common Rademacher sign $\varepsilon$,
$\varphi_\varepsilon(d\theta)=\cos(d\theta)$ and $\varphi_\varepsilon(\theta)^d=\cos^d\theta$, so the one-variable problem
already attains $C_d$; the same law is exchangeable and has identical marginals.
\end{proof}

\section{The optimizer and a continuous interpolation}
\begin{proof}[Proof of Theorem~\ref{thm:unique}]
Since $\sin d\theta=\sin\theta\,U_{d-1}(\cos\theta)$,
\[
g_d'(\theta)=d\sin\theta\,[U_{d-1}(\cos\theta)-\cos^{d-1}\theta].
\]
Using
\[
\frac{U_{d-1}(x)}{x^{d-1}}=2^{d-1}\prod_{k=1}^{\lfloor(d-1)/2\rfloor}
\left(1-\frac{\cos^2(k\pi/d)}{x^2}\right),\qquad x>0,
\]
(with the usual central factor when $d$ is even), the ratio is strictly increasing on
$(\cos(\pi/d),1]$, from $0$ to $d>1$. Hence it crosses $1$ exactly once, proving uniqueness.
The assertion about $|T_d-x^d|$ follows from Corollary~\ref{thm:main} and parity.

Strict growth of $C_d$ follows from Proposition~\ref{prop:continuous} below. Finally
$C_d\ge1+\cos^d(\pi/d)\to2$ and $C_d\le2$.
\end{proof}

For real $p>2$ define
\[
C(p)=\max_{0\le\theta\le\pi/p}\{\cos^p\theta-\cos(p\theta)\}
=\max_{0\le s\le\pi}\{\cos^p(s/p)-\cos s\}.
\]
Thus $C(d)=C_d$.

\begin{proposition}[Selected continuous consequences]\label{prop:continuous}
For $p>2$:
\begin{enumerate}
\item $C(p)$ is strictly increasing in $p$;
\item every maximizer $\theta_p$ is interior and satisfies
\[
\frac{\pi}{p+1}<\theta_p<\frac\pi p,
\qquad
\sin(p\theta_p)=\sin\theta_p\,\cos^{p-1}\theta_p;
\]
\item with $a_p=\pi/(p+1)$,
\[
\theta_p=a_p+\frac{\pi^3}{2p^3}+O(p^{-4}),\qquad
C(p)-\{\cos^pa_p+\cos a_p\}=\frac{\pi^6}{8p^4}+O(p^{-5}).
\]
\end{enumerate}
\end{proposition}
\begin{proof}
For fixed $s\in(0,\pi]$,
\[
\frac{\partial}{\partial p}\left[p\log\cos(s/p)\right]
=\log\cos z+z\tan z>0,\qquad z=s/p\in(0,\pi/2),
\]
because the derivative of $\log\cos z+z\tan z$ is $z\sec^2z>0$ and the expression vanishes at $z=0$.
Thus $\cos^p(s/p)$, and hence its maximum after subtracting $\cos s$, increases strictly with $p$.

At an endpoint $\theta=\pi/p$ the derivative is negative, whereas the objective is positive for small $\theta>0$;
therefore a maximizer is interior and satisfies the displayed stationary equation. Put
$\delta=\pi-p\theta_p$. Then
$\sin\delta=\sin\theta_p\cos^{p-1}\theta_p<\sin\theta_p$. Since
$0<\theta_p<\pi/p<\pi/2$ and $\delta<\pi-\theta_p$, this forces $\delta<\theta_p$, hence
$\theta_p>\pi/(p+1)$.

For the last assertion, expand the stationarity equation around $a_p$. One obtains
$f_p'(a_p)=\pi^3/(2p)+O(p^{-2})$ and $f_p''(a_p)=-p^2+O(p)$ for
$f_p(\theta)=\cos^p\theta-\cos(p\theta)$, hence
$\theta_p-a_p=-f_p'(a_p)/f_p''(a_p)+O(p^{-4})=\pi^3/(2p^3)+O(p^{-4})$.
Taylor expansion at the maximizer then gives the value correction $\pi^6/(8p^4)+O(p^{-5})$.
\end{proof}

\begin{remark}
The full continuous theory also gives a Gamma/Fourier realization and a two-parameter cosine/Chebyshev master
surface. Those results are deliberately not imported here: Proposition~\ref{prop:continuous} contains only the pieces
that sharpen the original sequence $C_d$.
\end{remark}

\section{Algebraic description and small dimensions}
\begin{corollary}\label{cor:small}
$C_2=1$, $C_3=2/\sqrt3$, $C_4=9/7$.
\end{corollary}
\begin{proof}
For $d=2$, $g_2(\theta)=\sin^2\theta$. For $d=3$, $x^3-T_3(x)=3x(1-x^2)$, maximal at $x=1/\sqrt3$. For $d=4$,
with $y=\cos^2\theta$, $g_4=-7y^2+8y-1$, maximal at $y=4/7$.
\end{proof}

\begin{proposition}\label{prop:alg}
For $d\ge3$, $c_d^*=\cos\theta_d^*$ is algebraic and
\[
\boxed{C_d=U_{d-2}(c_d^*).}
\]
Moreover $y=(c_d^*)^2$ satisfies an integer polynomial of degree $\lfloor(d-1)/2\rfloor$ with leading coefficient
$2^{d-1}-1$. In particular
\[
C_5=\tfrac4{15}(2\sqrt{21}-3)\sqrt{\tfrac{6+\sqrt{21}}{15}},\qquad
C_6=\frac{225+140\sqrt{70}}{961}.
\]
\end{proposition}
\begin{proof}
The critical equation is $U_{d-1}(c)=c^{d-1}$. For odd $d$ it is even in $c$; for even $d$ both sides are odd and we
divide by $c$, giving the stated polynomial in $y=c^2$. Finally
$T_d(c)=cU_{d-1}(c)-U_{d-2}(c)=c^d-U_{d-2}(c)$ at a critical point, so
$C_d=c^d-T_d(c)=U_{d-2}(c)$.
\end{proof}

\begin{remark}[The tangent-square critical polynomial]\label{rem:criticalpoly}
Writing $x=\tan^2\theta$ in the critical equation $\sin(d\theta)=\sin\theta\cos^{d-1}\theta$ gives
\[
 A_d(x):=\sum_{k=0}^{\lfloor(d-1)/2\rfloor}(-1)^k\binom d{2k+1}x^k-1=0.
\]
Thus $A_d$ is not introduced ad hoc: its distinguished positive root is exactly
$x_d^*=\tan^2\theta_d^*$, where $\theta_d^*$ is the optimizer defining $C_d$.
The arithmetic of this family---irreducibility, local Newton polygons, ramification and Galois groups---is developed
in \cite{MartinevskiCritical}. In particular, that work shows that the algebraic problem has substantial structure
well beyond the small-dimensional formulas recorded here.
\end{remark}

\section{The real case}
\begin{theorem}\label{thm:real}
For real random variables with $|X_j|\le1$,
\[
\sup\Bigl|\mathbb E\prod_jX_j-\prod_j\mathbb EX_j\Bigr|=1+\Bigl(1-\frac2d\Bigr)^d,
\]
attained when exactly one coordinate equals $-1$, at a uniformly random position, and the others equal $1$. The value
increases to $1+e^{-2}\approx1.1353$, which is strictly below $C_d$ for every $d\ge3$.
\end{theorem}
\begin{proof}
Randomizing each $X_j$ to $\pm1$ with conditional mean $X_j$ reduces to $U_j\in\{\pm1\}$. Put $A=\mathbb E\prod U_j$,
$a_j=\mathbb EU_j$, $B=\prod a_j$. If $AB\ge0$ then $|A-B|\le1$. Otherwise flipping $U_1$ changes $(A,B)$ to $(-A,-B)$, so
we may assume $B>0>A$; flipping negative $a_j$ in pairs we may assume all $a_j>0$. With $q_j=\mathbb P(U_j=-1)<\frac12$,
$N=\#\{j:U_j=-1\}$ and $\pi=\mathbb P(N\text{ even})$ we have $A=2\pi-1$ and $\sum_jq_j=\mathbb EN\ge1-\pi$. By AM--GM,
$B=\prod(1-2q_j)\le(1-\frac2d\sum q_j)^d\le(1-\frac{2(1-\pi)}d)^d$, so $B-A\le1+\psi(\pi)$ with
$\psi(\pi)=(1-2(1-\pi)/d)^d-2\pi$ non-increasing on $[0,\frac12]$; hence $B-A\le1+(1-2/d)^d$. The extremal law has $A=-1$,
$a_j=1-2/d$.
\end{proof}

\begin{remark}
For nonnegative factors a classical lemma of Ibragimov \cite{Ibragimov1962} controls the same functional under strong
mixing. Deviations from the piling-up lemma of linear cryptanalysis \cite{Matsui1994} under dependence were studied by
Kukorelly \cite{Kukorelly1999}; Theorem~\ref{thm:real} gives the sharp worst case for $\pm1$-valued variables.
\end{remark}

\section{Gaussian laws}
\begin{proposition}\label{prop:gauss-all}
If $X=(X_1,\dots,X_d)$ is Gaussian, then $D_d(X)\le1$, with strict inequality when the joint covariance
matrix is nonsingular. The value $1$ is attained (as a supremum over $t$) by singular Gaussian laws, e.g.\ $X_1=Z$,
$X_2=-Z$ with $Z\sim N(0,\sigma^2)$, for which $D_2=\sup_t(1-e^{-t^2\sigma^2})=1$.
\end{proposition}
\begin{proof}
With mean $\mu$, joint covariance $\Sigma$ and block covariances $\Sigma_{jj}$, after removing the common factor
$e^{i\langle t,\mu\rangle}$ both terms are real numbers in $(0,1]$:
$|\Delta_X(t)|=|e^{-t^\top\Sigma t/2}-e^{-\sum_jt_j^\top\Sigma_{jj}t_j/2}|<1$. If $\Sigma$ is nonsingular both
exponents tend to $+\infty$ together as $\|t\|\to\infty$ and stay bounded together on bounded sets, so the
supremum is attained and is strictly below one.
\end{proof}

For two blocks the defect is explicit. Let
$h(\lambda)=|\lambda-1|\lambda^{-\lambda/(\lambda-1)}$ for $\lambda>0$, $\lambda\ne1$, and $h(1)=0$.

\begin{proposition}[explicit form of the supremum in {\cite[Thm.~1(4)]{Daniusis2025}}]\label{prop:gauss-two}
Let $(X,Y)$ be jointly Gaussian with nonsingular block covariances and largest canonical correlation
$\rho_1\in[0,1)$. Then
\[
D_2(X,Y)=h(1-\rho_1)=\rho_1(1-\rho_1)^{(1-\rho_1)/\rho_1}\qquad(\rho_1>0),
\]
and $D_2=0$ if $\rho_1=0$. In particular $D_2\sim\rho_1/e$ as $\rho_1\downarrow0$ and $D_2\to1$ as
$\rho_1\uparrow1$.
\end{proposition}
\begin{proof}
Blockwise whitening does not change $D_2$ and reduces the joint covariance to
$R=\bigl(\begin{smallmatrix}I&C\\C^\top&I\end{smallmatrix}\bigr)$, whose eigenvalues lie in $[1-\rho_1,1+\rho_1]$
with both endpoints attained. Writing $t=ru$ with $\|u\|=1$, $\lambda=u^\top Ru$ and $x=r^2/2$,
$|\Delta(t)|=|e^{-\lambda x}-e^{-x}|$. For fixed $\lambda\ne1$ the maximum over $x\ge0$ is attained at
$x=\log\lambda/(\lambda-1)$ and equals $h(\lambda)$. One checks $h(\lambda)=h(1/\lambda)$, and $h$ increases on
$[1,\infty)$ because $e^{-x}-e^{-\lambda x}$ increases in $\lambda$ for each $x>0$. Since
$(1-\rho_1)(1+\rho_1)<1$, i.e.\ $1/(1-\rho_1)>1+\rho_1$, the maximum over $\lambda$ is $h(1-\rho_1)$.
\end{proof}

\begin{remark}
The same argument gives, for any two Gaussian laws with a common mean and positive definite covariances $A,B$,
$\sup_t|\varphi_{N(\mu,A)}(t)-\varphi_{N(\mu,B)}(t)|=\max\{h(\lambda_{\min}),h(\lambda_{\max})\}$, where
$\lambda_{\min},\lambda_{\max}$ are the extreme generalized eigenvalues of $Av=\lambda Bv$; equivalently it is an
increasing function of the Thompson distance $\|\log(B^{-1/2}AB^{-1/2})\|_{\rm op}$.
\end{remark}



\section{Large-\texorpdfstring{$d$}{d} behaviour}\label{sec:larged}
Let
\[
\widetilde C_d=\cos\frac{\pi}{d+1}+\cos^d\frac{\pi}{d+1}.
\]
The elementary point $\theta=\pi/(d+1)$ already gives $\widetilde C_d\le C_d$. We first record a
uniform finite-$d$ estimate and then the sharper asymptotic correction.

\begin{proposition}[Global approximant]\label{prop:globalapprox}
For every $d\ge4$,
\[
0\le C_d-\widetilde C_d\le\frac{\pi^6}{4d^4}.
\]
Moreover, with $\delta_d^*=\pi-d\theta_d^*$,
\[
\left|\delta_d^*-\frac{\pi}{d+1}\right|\le\frac{\pi^3}{2d^2}.
\]
\end{proposition}
\begin{proof}
For $\delta\in[0,\pi]$ put
\[
x(\delta)=\frac{\pi-\delta}{d},\qquad G(\delta)=\cos^d x(\delta)+\cos\delta.
\]
Then $C_d=\max_{[0,\pi]}G$. On $[\pi/2,\pi]$, $G\le1<1+\cos^d(\pi/d)=G(0)$, while
\[
G'(\delta)=\sin x\,\cos^{d-1}x-\sin\delta.
\]
Thus every maximizer $\delta_*\in(0,\pi/2)$ satisfies
$\sin\delta_*=\sin x_*\cos^{d-1}x_*$, where $x_*=x(\delta_*)\le\pi/d\le\pi/4$.

From $\sin\delta_*\le x_*\le\pi/d$ and $\sin\delta\ge2\delta/\pi$ on $[0,\pi/2]$,
$\delta_*\le\pi^2/(2d)$. Also $\sin\delta_*\ge\delta_*-\delta_*^3/6$ gives
\[
\delta_*\le\frac{\pi}{d+1}+\frac{\pi^6}{48d^3}.
\]
For the lower bound, Bernoulli's inequality and $x_*\le\pi/4$ give
$\cos^{d-1}x_*\ge1-(d-1)x_*^2/2\ge0$, and hence
\[
\delta_*\ge\sin\delta_*
\ge\left(x_*-\frac{x_*^3}{6}\right)
      \left(1-\frac{(d-1)x_*^2}{2}\right)
\ge x_*-\frac d2x_*^3
\ge\frac{\pi-\delta_*}{d}-\frac{\pi^3}{2d^2}.
\]
Therefore
\[
\left|\delta_*-\frac{\pi}{d+1}\right|\le\frac{\pi^3}{2d^2}.
\]
Finally
\[
G''(\delta)=-\frac1d\left[\cos^d x-(d-1)\sin^2x\cos^{d-2}x\right]-\cos\delta,
\]
and $|G''|\le2$ for $d\ge4$. Taylor's formula at $\delta_*$ therefore yields
\[
G\!\left(\frac{\pi}{d+1}\right)
\ge C_d-\left(\delta_*-\frac{\pi}{d+1}\right)^2
\ge C_d-\frac{\pi^6}{4d^4}.
\]
Since $G(\pi/(d+1))=\widetilde C_d\le C_d$, the proof is complete.
\end{proof}

\begin{proposition}[Sharp asymptotic correction]\label{prop:approx}
As $d\to\infty$,
\[
\theta_d^*=\frac{\pi}{d+1}+\frac{\pi^3}{2d^3}+O(d^{-4}),
\qquad
\boxed{C_d-\widetilde C_d=\frac{\pi^6}{8d^4}+O(d^{-5}).}
\]
Thus Proposition~\ref{prop:globalapprox} is sharp in order.
\end{proposition}
\begin{proof}
This is Proposition~\ref{prop:continuous}(3) with $p=d$.
\end{proof}

\begin{corollary}[Four-term expansion]\label{cor:asymp}
As $d\to\infty$,
\[
\begin{aligned}
C_d={}&2-\frac{\pi^2}{2d}
+\frac{\pi^2(\pi^2+4)}{8d^2}
-\frac{\pi^2(\pi^4+28\pi^2+24)}{48d^3}\\
&+\frac{\pi^2(\pi^6+112\pi^4+624\pi^2+192)}{384d^4}
+O(d^{-5}).
\end{aligned}
\]
Consequently
\[
C_{d+1}-C_d=\frac{\pi^2}{2d^2}-\frac{\pi^2(\pi^2+6)}{4d^3}+O(d^{-4}),
\]
and $C_d$ is strictly concave for all sufficiently large $d$.
\end{corollary}

\begin{remark}
Numerically the discrete concavity appears to hold already from $d=2$; a proof for every $d$ remains open.
\end{remark}

\section{Remarks}
\paragraph{Operator and dynamical forms.} Since $C_d=\|T_d-x^d\|_{\infty,[-1,1]}$ and the maximizers are $\pm c_d^*$
(Theorem~\ref{thm:unique}), the spectral theorem gives $\|T_d(A)-A^d\|\le C_d$ for every self-adjoint contraction $A$,
with equality for $A=c_d^*I$. This applies to reversible Markov operators on $L^2(\pi)$, where $T_d(P)$ describes
discrete-wave (quantum-walk) dynamics and $P^d$ the $d$-step random walk \cite{AM2024}; equality holds for the symmetric
two-state chain with eigenvalues $\{1,c_d^*\}$. Likewise $C_d=\sup_{|z|=1}|q(z^d)-q(z)^d|$ for $q(z)=(z+z^{-1})/2$, since
$q(z^d)=T_d(q(z))$ \cite{Rivlin1990}. These are restatements of Corollary~\ref{thm:main} and Theorem~\ref{thm:unique}, not
independent results.

\paragraph{Quantum parity.} For commuting local unitaries $V_j=\mathrm{diag}(e^{i\theta},e^{-i\theta})$ and the
Greenberger--Horne--Zeilinger state \cite{Mermin1990}, $\langle\otimes_jV_j\rangle=\cos d\theta$ and
$\langle V_j\rangle=\cos\theta$, so the factorization defect in this single state is $|\cos d\theta-\cos^d\theta|$ and equals
$C_d$ at $\theta=\theta_d^*$. For Hermitian observables with spectrum $\{\pm1\}$ the joint spectral law is
$\pm1$-valued, and the relevant sharp constant is the real one, $1+(1-2/d)^d$ (Theorem~\ref{thm:real}).

\paragraph{Normalization and estimation.} The ratio $D_d(X)/C_d$ lies in $[0,1]$ and vanishes exactly under mutual
independence; for $d=2$ it is the UFDM \cite{Daniusis2025}. It should be used with care: it is not invariant under
adding an independent coordinate (which leaves $D_d$ unchanged but replaces $C_d$ by $C_{d+1}$), it is not invariant under
monotone transformations of the marginals, and Gaussian laws never exceed $1/C_d$. Moreover, the supremum over all
frequencies cannot be estimated consistently from data: the empirical characteristic function is not uniformly
consistent on all of $\mathbb R^d$ \cite{Csorgo1981,CsorgoTotik1983}, and the uniform estimator of \cite{Daniusis2025}
restricts frequencies accordingly. In fact, if $N$ is even and, for each $j$, the observed values $x_{1j},\dots,x_{Nj}$
are linearly independent over $\mathbb Q$ (almost surely for i.i.d.\ samples with absolutely continuous marginals), then
the unrestricted empirical supremum equals $C_d$ for every such data set, by Kronecker's theorem and
Corollary~\ref{thm:main}.

\paragraph{Mixing.} Corollary~\ref{thm:main} is an unconditional ceiling for multi-factor covariance inequalities under
strong mixing, such as the Volkonskii--Rozanov lemma \cite{VolkonskiiRozanov1959}; sharp mixing constants are treated
separately.

\section{Open problems}
\begin{enumerate}
\item Classify all extremizers of the profile/universal theorem for $d\ge3$ and prove a quantitative stability theorem.
The $d=2$ extremal set is larger.
\item Prove irreducibility of the critical polynomial $A_d$ for every $d$. A companion paper proves broad infinite families, local factorization theorems and large Galois-group results, but the all-$d$ conjecture remains open \cite{MartinevskiCritical}.
\item Prove or refute global discrete concavity $C_{d+2}-2C_{d+1}+C_d<0$ for every $d\ge2$.
\item Extend the structural restrictions beyond the classes already settled in \cite{MartinevskiMarginals}: finite phase alphabets, further fixed-marginal families, and broader Gaussian/elliptical classes.
\end{enumerate}

\section*{Supplementary material}
A table of $C_d$, $\theta_d^*$ and $c_d^*$ for $2\le d\le100$, computation scripts, and the previous independent
phase-shift/angular proof of Corollary~\ref{thm:main} accompany this manuscript. The older proof is retained as an
independent verification route, not as the primary proof.

\section*{Historical note}
The uniform characteristic-function factorization defect and the question of its extremal constants, including the
limit $2$, go back to an unpublished student research project of the author (1996), originally motivated in part by
cryptographic applications. The exact finite-$d$ formula, fixed-mean profile, proofs, continuous refinements and the
critical-polynomial formulation in this paper were obtained in 2026.

\section*{Author information and AI assistance}
\noindent
\textbf{Author.} Alexandr Martinevski, Independent Researcher, Lithuania. ORCID:
\href{https://orcid.org/0009-0000-4230-4414}{0009-0000-4230-4414}.
Email: \href{mailto:alex.martinevski@gmail.com}{alex.martinevski@gmail.com}.
Manuscript date: September 2026.

\medskip
\noindent
\textbf{AI assistance disclosure.} AI systems used in the research and manuscript workflow included
\href{https://openai.com/index/gpt-5-6/}{OpenAI GPT-5.6 Sol} and
\href{https://www.anthropic.com/claude-opus-5-5}{Anthropic Claude Opus 5.5}.
They were used for mathematical exploration, proof development and critique, symbolic and numerical verification,
literature search, and manuscript drafting and editing. AI systems are not authors. The author reviewed the
mathematical statements and proofs and takes full responsibility for the content.

\begin{thebibliography}{99}\small
\bibitem{AM2024} S. Apers and L. Miclo, Quantum walks, the discrete wave equation and Chebyshev polynomials, arXiv:2402.07809 (2024).
\bibitem{BKRS2019} B. B\"ottcher, M. Keller-Ressel and R.~L. Schilling, Distance multivariance: new dependence measures for random vectors, \emph{Ann. Statist.} 47 (2019) 2757--2789.
\bibitem{Csorgo1981} S. Cs\"org\H{o}, Limit behaviour of the empirical characteristic function, \emph{Ann. Probab.} 9 (1981) 130--144.
\bibitem{Csorgo1985} S. Cs\"org\H{o}, Testing for independence by the empirical characteristic function, \emph{J. Multivariate Anal.} 16 (1985) 290--299.
\bibitem{CsorgoTotik1983} S. Cs\"org\H{o} and V. Totik, On how long interval is the empirical characteristic function uniformly consistent?, \emph{Acta Sci. Math. (Szeged)} 45 (1983) 141--149.
\bibitem{Daniusis2025} P. Daniu\v{s}is, S. Juneja, L. Kuzma and V. Marcinkevi\v{c}ius, Measuring statistical dependence via characteristic function IPM, \emph{Entropy} 27(12) (2025) 1254.
\bibitem{EHSV2010} N. Ebrahimi, G.~G. Hamedani, E.~S. Soofi and H. Volkmer, A class of models for uncorrelated random variables, \emph{J. Multivariate Anal.} 101 (2010) 1859--1871.
\bibitem{Hamedani2013} G.~G. Hamedani, Sub-independence: an expository perspective, \emph{Comm. Statist. Theory Methods} 42 (2013) 3615--3638.
\bibitem{Ibragimov1962} I.~A. Ibragimov, Some limit theorems for stationary processes, \emph{Theory Probab. Appl.} 7 (1962) 349--382.
\bibitem{Kankainen1995} A. Kankainen, \emph{Consistent Testing of Total Independence Based on the Empirical Characteristic Function}, PhD thesis, University of Jyv\"askyl\"a, 1995.
\bibitem{Kukorelly1999} Z. Kukorelly, The piling-up lemma and dependent random variables, in \emph{Cryptography and Coding}, LNCS 1746, Springer, 1999, 186--190.
\bibitem{Matsui1994} M. Matsui, Linear cryptanalysis method for DES cipher, \emph{EUROCRYPT '93}, LNCS 765 (1994) 386--397.
\bibitem{Mermin1990} N.~D. Mermin, Extreme quantum entanglement in a superposition of macroscopically distinct states, \emph{Phys. Rev. Lett.} 65 (1990) 1838--1840.
\bibitem{MartinevskiMarginals} A. Martinevski, \emph{Two-point extremals for characteristic-function dependence with fixed marginals}, companion manuscript (2026).
\bibitem{MartinevskiCritical} A. Martinevski, \emph{Chebyshev--Lucas critical polynomials: local blocks, irreducibility and Galois groups}, companion manuscript (2026).
\bibitem{Pfister2018} N. Pfister, P. B\"uhlmann, B. Sch\"olkopf and J. Peters, Kernel-based tests for joint independence, \emph{J. R. Stat. Soc. Ser. B} 80 (2018) 5--31.
\bibitem{Rivlin1990} T.~J. Rivlin, \emph{Chebyshev Polynomials}, 2nd ed., Wiley, 1990.
\bibitem{Schennach2019} S.~M. Schennach, Convolution without independence, \emph{J. Econometrics} 211(1) (2019) 308--318.
\bibitem{VolkonskiiRozanov1959} V.~A. Volkonskii and Yu.~A. Rozanov, Some limit theorems for random functions I, \emph{Theory Probab. Appl.} 4 (1959) 178--197.
\end{thebibliography}
\end{document}
